% =====================================================================
%  Algorithmic Game Theory — Progetto d'esame (documento di lavoro)
%  Compilazione:  powershell -File compila.ps1 project   ->  project.pdf
%
%  Si aggiorna dopo ogni lezione: nuove idee nella "banca delle idee",
%  risposte alle domande aperte, decisioni prese.
% =====================================================================
\documentclass[11pt, a4paper, oneside, openany]{book}
% =====================================================================
%  preambolo.tex — pacchetti, stile e ambienti comuni a tutto il documento
% =====================================================================

% ---------- Lingua e codifica ----------
\usepackage[T1]{fontenc}
\usepackage[utf8]{inputenc}
\usepackage[italian]{babel}

% ---------- Font ----------
% Times-like (newtx): tratto pieno e scuro, larghezza compatta.
% Anche la matematica è coerente col testo.
\usepackage{newtxtext}
\usepackage{newtxmath}
\usepackage[scaled=0.88]{beramono}   % monospace leggibile per il codice
\usepackage{microtype}

% ---------- Pagina ----------
\usepackage[a4paper, top=2.5cm, bottom=2.5cm, left=2.5cm, right=2.5cm, headheight=14pt]{geometry}
\setlength{\parindent}{0pt}
\setlength{\parskip}{0.55em}
\linespread{1.08}

% ---------- Colori ----------
\usepackage[dvipsnames]{xcolor}
\definecolor{unical}{HTML}{B00000}     % rosso delle slide
\definecolor{blu}{HTML}{0B3A91}        % blu delle slide di sezione
\definecolor{verde}{HTML}{1E7B34}
\definecolor{arancio}{HTML}{C0620A}
\definecolor{grigio}{HTML}{4A4A4A}
\definecolor{sfondo}{HTML}{F5F6F8}

% ---------- Matematica ----------
\usepackage{amsmath}
\usepackage{mathtools}

% ---------- Tabelle, liste, figure ----------
\usepackage{booktabs}
\usepackage{array}
\usepackage{multirow}
\usepackage{tabularx}
\usepackage{enumitem}
\setlist{itemsep=0.2em, topsep=0.3em}
\setlist[itemize,1]{label=\textcolor{unical}{\small$\blacktriangleright$}}
\usepackage{graphicx}
\usepackage{float}
\usepackage{caption}
\usepackage{needspace}   % \needspace{n\baselineskip}: va a capo pagina se manca spazio
\captionsetup{font=small, labelfont={bf,color=blu}}

\usepackage{tikz}
\usetikzlibrary{arrows.meta, positioning, trees, calc, shapes.geometric}

% ---------- Titoli ----------
\usepackage{titlesec}
% Capitolo = una lezione
\titleformat{\chapter}[display]
  {\normalfont\sffamily\bfseries\color{blu}}
  {\large\MakeUppercase{\chaptertitlename}\ \thechapter}
  {0.3em}
  {\Huge}
  [\vspace{0.3em}{\color{unical}\titlerule[1.2pt]}]
\titlespacing*{\chapter}{0pt}{0pt}{1.5em}
\titleformat{\section}
  {\normalfont\sffamily\Large\bfseries\color{blu}}{\thesection}{0.8em}{}
\titleformat{\subsection}
  {\normalfont\sffamily\large\bfseries\color{black}}{\thesubsection}{0.8em}{}
\titleformat{\subsubsection}
  {\normalfont\sffamily\normalsize\bfseries\color{grigio}}{\thesubsubsection}{0.8em}{}
\addto\captionsitalian{\renewcommand{\chaptername}{Lezione}}

% ---------- Intestazioni ----------
\usepackage{fancyhdr}
\pagestyle{fancy}
\fancyhf{}
\fancyhead[L]{\small\sffamily\color{grigio}Algorithmic Game Theory}
\fancyhead[R]{\small\sffamily\color{grigio}\leftmark}
\fancyfoot[C]{\small\sffamily\thepage}
\renewcommand{\headrulewidth}{0.4pt}
\renewcommand{\chaptermark}[1]{\markboth{Lezione \thechapter\ -- #1}{}}
\fancypagestyle{plain}{\fancyhf{}\fancyfoot[C]{\small\sffamily\thepage}\renewcommand{\headrulewidth}{0pt}}

% ---------- Box (definizioni, esempi, ...) ----------
\usepackage[most]{tcolorbox}
\tcbset{
  appunti/.style={
    enhanced, breakable, sharp corners,
    boxrule=0pt, leftrule=3pt,
    left=8pt, right=8pt, top=5pt, bottom=5pt,
    before skip=0.8em, after skip=0.8em,
    fonttitle=\sffamily\bfseries, coltitle=black,
    attach title to upper, after title={.\ },
  }
}
\newtcolorbox[auto counter, number within=chapter]{definizione}[1][]{
  appunti, colback=blu!5, colframe=blu,
  title={\textcolor{blu}{Definizione~\thetcbcounter}\ifx\relax#1\relax\else\ (#1)\fi}}
\newtcolorbox[use counter from=definizione]{teorema}[1][]{
  appunti, colback=unical!5, colframe=unical,
  title={\textcolor{unical}{Teorema~\thetcbcounter}\ifx\relax#1\relax\else\ (#1)\fi}}
\newtcolorbox{esempio}[1][]{
  appunti, colback=verde!5, colframe=verde,
  title={\textcolor{verde}{Esempio}\ifx\relax#1\relax\else\ (#1)\fi}}
\newtcolorbox{osservazione}[1][]{
  appunti, colback=sfondo, colframe=grigio,
  title={\textcolor{grigio}{Osservazione}\ifx\relax#1\relax\else\ (#1)\fi}}
% Punti delicati, errori tipici, imprecisioni nelle slide
\newtcolorbox{attenzione}[1][]{
  appunti, colback=arancio!7, colframe=arancio,
  title={\textcolor{arancio}{Attenzione}\ifx\relax#1\relax\else\ (#1)\fi}}
% Contenuti aggiunti negli appunti, non presenti nelle slide
\newtcolorbox{extra}[1][]{
  appunti, colback=white, colframe=grigio!60, leftrule=1.5pt,
  borderline west={0.4pt}{4pt}{grigio!60},
  title={\textcolor{grigio}{Approfondimento}\ifx\relax#1\relax\else\ (#1)\fi}}
% Riepilogo di fine lezione
\newtcolorbox{riepilogo}{
  enhanced, breakable, sharp corners, colback=blu!4, colframe=blu, boxrule=0.8pt,
  fonttitle=\sffamily\bfseries, title={Riepilogo della lezione}}

% ---------- Teoria vs Laboratorio ----------
% \parteteoria e \partelab aprono le due metà di una lezione con un'etichetta colorata.
\newcommand{\etichetta}[2]{%
  \tikz[baseline=(t.base)]\node[fill=#1, text=white, font=\sffamily\bfseries\small,
    inner xsep=6pt, inner ysep=3pt](t){#2};}
\newcommand{\parteteoria}{%
  \par\medskip\noindent\etichetta{blu}{TEORIA}\par\smallskip}
\newcommand{\partelab}[1]{%
  \clearpage
  \section{Laboratorio: #1}
  \noindent\etichetta{unical}{LAB}\par\smallskip}
\newcommand{\fonte}[1]{\hfill{\small\sffamily\color{grigio}[#1]}}

% ---------- Codice ----------
\usepackage{listings}
\lstdefinestyle{python}{
  language=Python,
  basicstyle=\ttfamily\small,
  keywordstyle=\color{blu}\bfseries,
  commentstyle=\color{verde}\itshape,
  stringstyle=\color{unical},
  numberstyle=\tiny\color{grigio}, numbers=left, numbersep=8pt,
  backgroundcolor=\color{sfondo},
  frame=leftline, framerule=2pt, rulecolor=\color{blu},
  xleftmargin=1.8em, framexleftmargin=1.8em,
  showstringspaces=false, breaklines=true, tabsize=4,
  columns=fullflexible, keepspaces=true,
  morekeywords={True,False,None},
  literate={à}{{\`a}}1 {è}{{\`e}}1 {é}{{\'e}}1 {ì}{{\`i}}1 {ò}{{\`o}}1 {ù}{{\`u}}1
           {−}{{-}}1 {∞}{{$\infty$}}1 {×}{{$\times$}}1,
}
\lstdefinestyle{pseudo}{
  style=python, numbers=none, mathescape=true,
  morekeywords={def,if,else,for,each,return,break,in,and,or,not},
}
\lstset{style=python}

% ---------- Tris (tic-tac-toe) ----------
% \tris{X,O,e, e,X,e, O,e,e}: 9 celle riga per riga, e = vuota.
\newcommand{\trisX}{\textbf{X}}
\newcommand{\trisO}{\textcolor{unical}{\textbf{O}}}
\newcommand{\trise}{}
\newcommand{\tris}[2][0.42]{%
  \begin{tikzpicture}[scale=#1, baseline=(current bounding box.center)]
    \draw[thick] (1,0)--(1,3) (2,0)--(2,3) (0,1)--(3,1) (0,2)--(3,2);
    \foreach \c [count=\k from 0] in {#2} {
      \pgfmathtruncatemacro\riga{\k/3}
      \pgfmathtruncatemacro\colonna{mod(\k,3)}
      \node[font=\sffamily] at (\colonna+0.5, 2.5-\riga) {\csname tris\c\endcsname};
    }
  \end{tikzpicture}}

% ---------- Utilità ----------
\newcommand{\en}[1]{\textit{#1}}            % termine inglese originale
\newcommand{\kw}[1]{\textbf{#1}}            % parola chiave
\newcommand{\R}{\mathbb{R}}
\newcommand{\N}{\mathbb{N}}
\newcommand{\Ev}{\mathbb{E}}
% Relazioni di preferenza
\newcommand{\pref}{\succsim}                % x pref y: x e' almeno buono quanto y
\newcommand{\npref}{\not\succsim}
\newcommand{\spref}{\succ}                  % preferenza stretta
\newcommand{\indif}{\sim}                   % indifferenza
\newcommand{\lott}{\mathcal{L}}             % insieme delle lotterie semplici
% Teoria dei giochi
\newcommand{\BR}{\mathrm{BR}}               % best response
\newcommand{\uno}{\mathbf{1}}               % vettore di tutti 1
\DeclareMathOperator{\supp}{supp}           % supporto di una strategia mista
\DeclareMathOperator{\core}{Core}
% Tabella dei payoff di un gioco a due giocatori: celle "a, b"
\newcommand{\pay}[2]{$#1,\,#2$}
\DeclareMathOperator*{\argmax}{arg\,max}
\DeclareMathOperator*{\argmin}{arg\,min}

% ---------- Link (per ultimo) ----------
\usepackage[hidelinks]{hyperref}
\hypersetup{pdftitle={Algorithmic Game Theory -- Appunti}, pdflang={it}}
\usepackage{bookmark}


% In questo documento i capitoli non sono lezioni
\addto\captionsitalian{\renewcommand{\chaptername}{Capitolo}}
\renewcommand{\chaptermark}[1]{\markboth{Progetto -- #1}{}}

% Stato delle idee e delle domande
\newcommand{\stato}[2]{\etichetta{#1}{#2}}
\newcommand{\aperta}{\stato{arancio}{APERTA}}
\newcommand{\risolta}{\stato{verde}{RISOLTA}}
\newcommand{\sbloccata}[1]{\stato{blu}{DA LEZIONE #1}}

\begin{document}
\frontmatter

\hypersetup{pageanchor=false}
\begin{titlepage}
  \centering
  \vspace*{3cm}
  {\sffamily\bfseries\color{unical}\large UNIVERSITÀ DELLA CALABRIA\par}
  \vspace{0.3cm}
  {\sffamily\color{grigio}Anno accademico 2026/2027 -- Appello di gennaio\par}
  \vspace{2.5cm}
  {\color{blu}\rule{\linewidth}{1.2pt}\par}
  \vspace{0.6cm}
  {\sffamily\bfseries\Huge\color{blu} Progetto d'esame\par}
  \vspace{0.4cm}
  {\sffamily\Large Algorithmic Game Theory\par}
  \vspace{0.6cm}
  {\color{blu}\rule{\linewidth}{1.2pt}\par}
  \vspace{1.5cm}
  {\large Documento di lavoro: requisiti, metodo, idee, decisioni\par}
  \vfill
  {\small\color{grigio} Aggiornato fino alla Lezione 5 \quad·\quad Ultima compilazione: \today\par}
\end{titlepage}
\hypersetup{pageanchor=true}

\tableofcontents

\mainmatter
% =====================================================================
\chapter{Cosa sappiamo dell'esame}

L'unica informazione ufficiale finora (slide della Lezione~1):

\begin{definizione}[Esame]
L'esame consiste in un \kw{progetto} in cui si deve \kw{formalizzare} e \kw{implementare}
la soluzione di un problema usando gli strumenti dell'Algorithmic Game Theory.
\end{definizione}

Due verbi, due metà del lavoro:
\begin{itemize}
  \item \kw{formalizzare}: trasformare un problema a parole in un modello preciso
        (giocatori, azioni, utilità, concetto di soluzione);
  \item \kw{implementare}: scrivere codice che calcola quella soluzione, e mostrare che funziona.
\end{itemize}

\begin{attenzione}[cosa non sappiamo ancora]
Formato della consegna, lavoro singolo o di gruppo, chi sceglie il problema, scadenze e
criteri di valutazione non sono ancora stati comunicati. Le domande da fare al docente sono
raccolte nel Capitolo~\ref{cap:domande}. Finché non ci sono risposte, tutto ciò che segue è
una \emph{proposta} di metodo, non un requisito.
\end{attenzione}

% =====================================================================
\chapter{Come impostare un progetto AGT}
\label{cap:metodo}

Qualunque sia il tema, un progetto di AGT segue la stessa sequenza. Conviene usarla anche
come scaletta della relazione finale.

\section{Le sette domande della formalizzazione}

\begin{enumerate}[label=\textbf{\arabic*.}, leftmargin=2em]
  \item \kw{Problema}: qual è la situazione reale? Perché è interessante?
  \item \kw{Giocatori}: chi prende decisioni? Quanti sono? ($N = \{1,\dots,n\}$)
  \item \kw{Azioni / strategie}: cosa può fare ciascuno? ($A_i$, oppure mosse in uno stato)
  \item \kw{Informazione e tempo}: si muove a turni o simultaneamente? Ognuno sa tutto
        (informazione perfetta) o no? C'è casualità?
  \item \kw{Utilità}: cosa vuole ciascun giocatore? ($u_i$, oppure $v(S)$ per le coalizioni).
        È un gioco a somma zero?
  \item \kw{Concetto di soluzione}: cosa vogliamo calcolare? (valore minimax, equilibrio di
        Nash, valore di Shapley, matching stabile, vincitore di un'elezione, \ldots)
  \item \kw{Algoritmo}: come lo calcoliamo? Con che complessità? Si può fare in modo esatto
        o serve un'approssimazione o un'euristica?
\end{enumerate}

\begin{esempio}[il tris del Lab 1 formalizzato]
\par\smallskip\noindent
\begin{tabularx}{\linewidth}{@{}>{\bfseries}l X@{}}
Problema      & giocare a tris in modo ottimale. \\
Giocatori     & 2: MAX (X) e MIN (O). \\
Azioni        & piazzare il proprio segno in una cella vuota. \\
Informazione  & a turni alternati, informazione perfetta, niente caso. \\
Utilità       & $+1$ / $0$ / $-1$ per MAX; somma zero. \\
Soluzione     & valore minimax di ogni stato, e la mossa che lo realizza. \\
Algoritmo     & Minimax con potatura alpha-beta; l'albero completo ha $549\,946$ nodi. \\
\end{tabularx}
\end{esempio}

\section{Dall'implementazione alla relazione}

Dopo la formalizzazione:
\begin{enumerate}[resume*]
  \item \kw{Implementazione}: codice pulito, eseguibile, con esempi e test.
  \item \kw{Esperimenti}: misure che mostrano che l'algoritmo funziona e quanto costa
        (tempo, nodi visitati, qualità della soluzione) al crescere dell'istanza.
  \item \kw{Discussione}: limiti del modello (quali ipotesi sono irrealistiche?) e del metodo
        (dove non scala?), possibili estensioni.
\end{enumerate}

\begin{osservazione}[cosa rende un progetto convincente]
\begin{itemize}
  \item Un confronto: algoritmo base contro versione migliorata, oppure due concetti di soluzione
        sullo stesso problema. Esempio già visto: Minimax contro alpha-beta, $549\,946$ nodi
        contro $18\,297$.
  \item Numeri misurati, non solo affermati.
  \item Ipotesi del modello dette esplicitamente, con il motivo per cui sono ragionevoli.
\end{itemize}
\end{osservazione}

\section{Checklist prima della consegna}

\begin{itemize}
  \item[$\square$] Ogni elemento del modello (giocatori, azioni, utilità) è definito in modo preciso.
  \item[$\square$] Il concetto di soluzione è motivato: perché proprio questo?
  \item[$\square$] L'algoritmo è descritto (pseudocodice) e se ne discute la complessità.
  \item[$\square$] Il codice gira da zero seguendo le istruzioni.
  \item[$\square$] Ci sono esperimenti con tabelle o grafici e un commento ai risultati.
  \item[$\square$] Ci sono limiti ed estensioni.
\end{itemize}

% =====================================================================
\chapter{Banca delle idee}
\label{cap:idee}

Le idee sono ordinate per argomento del corso. Ogni lezione ``sblocca'' nuovi strumenti:
un'idea diventa praticabile quando il suo argomento è stato spiegato.

\section{Giochi combinatori (Lezione 1)}

Strumenti disponibili: albero di gioco, Minimax, alpha-beta, e le tecniche per giochi grandi
(profondità limitata, euristiche, \en{transposition table}, \en{move ordering}).

\subsection*{Idea 1: IA per Forza 4 \hfill\sbloccata{1}}
\begin{itemize}
  \item \kw{Modello}: identico al tris (2 giocatori, turni, informazione perfetta, somma zero),
        ma su una griglia $6\times7$ con gravità.
  \item \kw{Perché è interessante}: l'albero completo è troppo grande per Minimax puro, quindi
        servono tutte le tecniche della slide ``Oltre alpha-beta''.
  \item \kw{Esperimenti}: nodi visitati e tempo al variare della profondità; effetto di
        \en{move ordering} (colonne centrali prima) e della \en{transposition table};
        percentuale di vittorie contro un giocatore casuale e uno ``greedy''.
  \item \kw{Rischio}: basso. È un'estensione diretta del Lab~1.
\end{itemize}

\subsection*{Idea 2: tris generalizzato $(m,n,k)$ \hfill\sbloccata{1}}
\begin{itemize}
  \item \kw{Modello}: griglia $m\times n$, vince chi allinea $k$ segni (il tris è $(3,3,3)$,
        il \en{gomoku} è $(15,15,5)$).
  \item \kw{Perché è interessante}: permette di misurare come esplode l'albero di gioco al
        crescere di $m, n, k$, e da che punto la ricerca esatta smette di essere praticabile.
  \item \kw{Esperimenti}: per istanze piccole, valore esatto del gioco (vince il primo o si pareggia?);
        per istanze grandi, ricerca a profondità limitata con euristica.
  \item \kw{Rischio}: basso-medio. La parte esatta è facile; l'euristica richiede tentativi.
\end{itemize}

\subsection*{Idea 3: IA per Othello (Reversi) o Dama \hfill\sbloccata{1}}
\begin{itemize}
  \item \kw{Modello}: gioco combinatorio con regole di mossa più ricche.
  \item \kw{Perché è interessante}: la funzione di valutazione euristica diventa centrale
        (mobilità, angoli, stabilità dei pezzi in Othello).
  \item \kw{Rischio}: medio. Implementare correttamente le regole richiede tempo.
\end{itemize}

\begin{osservazione}[limite dei giochi combinatori come progetto]
Queste idee usano solo la prima parte del corso. Il resto del programma (Nash, Shapley,
voto, matching, meccanismi, divisione equa) offrirà probabilmente temi più vicini allo spirito
``formalizzare un problema reale''. Conviene tenerle come piano B e decidere più avanti.
\end{osservazione}

\section{Teoria dell'utilità e decisioni sotto incertezza (Lezione 2)}

Strumenti disponibili: relazioni di preferenza, utilità ordinale e cardinale (VNM),
lotterie e utilità attesa, atteggiamento verso il rischio, ricerca bayesiana
(massimizzazione dell'utilità attesa con posteriore aggiornato).

\subsection*{Idea 4: ricerca bayesiana per un problema di \en{search and rescue} \hfill\sbloccata{2}}
\begin{itemize}
  \item \kw{Modello}: un solo decisore contro l'incertezza. Un oggetto (relitto, persona
        dispersa, guasto in una rete) si trova in una cella sconosciuta con un prior non
        uniforme; ogni ispezione ha un costo e una probabilità di \en{falso negativo}.
  \item \kw{Perché è interessante}: generalizza direttamente il Lab~2 e lo avvicina al caso
        reale dei paper in \texttt{labs/02-battleship/papers/} (volo AF~447). Permette di
        introdurre cose
        che nel lab non ci sono: rilevazione imperfetta, costi di spostamento diversi per
        cella, vincolo di budget.
  \item \kw{Concetto di soluzione}: politica che massimizza l'utilità attesa; confronto tra
        la politica \en{greedy} (massima probabilità) e una che tiene conto del costo.
  \item \kw{Esperimenti}: turni o costo medio per trovare l'oggetto, al variare del prior,
        del tasso di falsi negativi e del budget; confronto con le euristiche del Lab~2.
  \item \kw{Rischio}: basso. Il codice del Lab~2 è già una base funzionante.
\end{itemize}

\subsection*{Idea 5: atteggiamento verso il rischio in un problema di allocazione \hfill\sbloccata{2}}
\begin{itemize}
  \item \kw{Modello}: più agenti con funzioni di Bernoulli diverse ($u(x)=x$, $u(x)=\sqrt{x}$,
        $u(x)=x^2$) devono scegliere tra portafogli o contratti rischiosi.
  \item \kw{Perché è interessante}: mostra in modo misurabile che la soluzione ottima dipende
        dalla \emph{curvatura} di $u$ e non dal valore atteso, e permette di calcolare
        equivalenti certi e premi per il rischio.
  \item \kw{Collegamento}: è il ponte naturale verso i giochi coalizionali (il fondo comune
        della lezione è già una coalizione) e verso il \en{mechanism design} (le assicurazioni).
  \item \kw{Rischio}: medio. Da solo è un po' esile: funziona meglio come \emph{componente}
        di un progetto più grande, una volta visti Nash o Shapley.
\end{itemize}

\section{Giochi non cooperativi (Lezioni 3 e 4)}

Strumenti disponibili: forma normale, dominanza e IESDS, best response, equilibri di Nash
puri e misti, condizione del supporto, minimax con programmazione lineare, giochi
potenziali e dinamica di best response, equilibri correlati con programmazione lineare,
fictitious play e regret matching. Codice già scritto e verificato nei laboratori 3 e 4.

\subsection*{Idea 6: giochi di sicurezza (pattugliamento) \hfill\sbloccata{3}}
\begin{itemize}
  \item \kw{Modello}: un difensore con $k$ risorse (pattuglie, guardie) deve proteggere $n$
        bersagli; un attaccante sceglie quale colpire. Ogni bersaglio ha un valore diverso
        per i due giocatori. È uno degli esempi di applicazione citati nella lezione 1.
  \item \kw{Perché è interessante}: è un problema reale (aeroporti, porti, parchi naturali)
        e mostra il valore della randomizzazione: un difensore prevedibile viene aggirato.
        Nel caso a somma zero il minimax del laboratorio 3 risolve direttamente il problema.
  \item \kw{Concetto di soluzione}: strategia minimax del difensore; confronto con la
        strategia ``ingenua'' che protegge sempre i bersagli più preziosi.
  \item \kw{Esperimenti}: perdita attesa al variare di $k$ e di $n$; tempo del programma
        lineare quando le strategie pure del difensore (sottoinsiemi di $k$ bersagli)
        esplodono combinatoriamente, il che porta naturalmente a rappresentazioni compatte.
  \item \kw{Rischio}: basso-medio. La base (minimax con \texttt{linprog}) c'è già; la parte
        non banale è la crescita del numero di strategie.
\end{itemize}

\subsection*{Idea 7: rilevamento di comunità come gioco potenziale \hfill\sbloccata{3}}
\begin{itemize}
  \item \kw{Modello}: il gioco del laboratorio 3, con nodi che scelgono un cluster. Si
        estende con utilità diverse (per esempio penalizzando i cluster troppo grandi, per
        evitare l'equilibrio banale ``tutti insieme'').
  \item \kw{Perché è interessante}: il laboratorio ha mostrato che l'equilibrio raggiunto
        dipende dal punto di partenza e che il massimo del potenziale è la soluzione
        sbagliata. Studiare quando la dinamica trova le comunità giuste è una domanda vera.
  \item \kw{Esperimenti}: modelli a blocchi stocastici con rapporto $p/q$ variabile;
        accuratezza (per esempio con l'indice NMI) e numero di turni; confronto con Louvain
        di \texttt{networkx}.
  \item \kw{Rischio}: basso. Il codice c'è e gira in pochi secondi su 400 nodi.
\end{itemize}

\subsection*{Idea 8: apprendimento negli equilibri \hfill\sbloccata{4}}
\begin{itemize}
  \item \kw{Modello}: giochi ripetuti; giocatori che usano fictitious play, regret matching o
        varianti, anche uno contro l'altro.
  \item \kw{Perché è interessante}: il laboratorio 4 ha mostrato che una variante del regret
        matching (quella delle slide) converge al punto sbagliato su un gioco asimmetrico. Un
        confronto sistematico (giochi a somma zero, coordinamento, giochi con equilibri
        correlati diversi da Nash) dice quando e a che cosa converge ogni algoritmo.
  \item \kw{Esperimenti}: distanza dall'equilibrio nel tempo; distribuzione congiunta del
        regret matching confrontata con il politopo degli equilibri correlati calcolato con
        programmazione lineare; giochi con cicli (Shapley $3\times3$, dove fictitious play
        non converge).
  \item \kw{Rischio}: basso. Molto ``sperimentale'', meno ``modellazione di un problema
        reale'': da chiarire se è nello spirito richiesto.
\end{itemize}

\begin{osservazione}[il semaforo come progetto]
Il calcolo degli equilibri correlati con l'obiettivo maximin ha restituito esattamente il
semaforo nell'attraversamento. Un progetto che progetta dispositivi di coordinamento
(incroci, accesso a un canale radio, turni di utilizzo di una risorsa) mettendo a confronto
benessere ed equità sarebbe una buona sintesi delle lezioni 3 e 4.
\end{osservazione}

\section{Giochi coalizionali (Lezione 5)}

Strumenti disponibili: giochi in funzione caratteristica, superadditività, struttura di
coalizione con programmazione dinamica, core, nucleolo (Kopelowitz), valore di Shapley.
Codice già scritto e verificato nel laboratorio 5, compresi la programmazione dinamica e il
nucleolo.

\subsection*{Idea 9: ripartizione equa di costi o ricavi \hfill\sbloccata{5}}
\begin{itemize}
  \item \kw{Modello}: un problema concreto di condivisione, per esempio aziende che
        condividono un impianto, utenti che si dividono la bolletta di un servizio comune,
        comuni che finanziano un'infrastruttura. La funzione caratteristica si ricava da un
        modello di costo (spesso analitico).
  \item \kw{Perché è interessante}: permette di confrontare sullo stesso problema tre
        concetti di soluzione (core, nucleolo, Shapley) e di misurare il conflitto tra
        stabilità ed equità, visto più volte nella lezione.
  \item \kw{Esperimenti}: il core è vuoto? Shapley vi appartiene? Quanto distano nucleolo e
        Shapley? Tempi di calcolo al crescere di $n$ e approssimazione di Shapley per
        campionamento di permutazioni quando $2^n$ diventa troppo.
  \item \kw{Rischio}: basso-medio. La parte delicata è costruire una funzione caratteristica
        credibile.
\end{itemize}

\subsection*{Idea 10: valore di Shapley per spiegare un modello (SHAP) \hfill\sbloccata{5}}
\begin{itemize}
  \item \kw{Modello}: le feature di un modello di machine learning sono i giocatori, $v(C)$
        è la qualità della predizione usando solo le feature in $C$. È l'applicazione
        citata nella lezione 1.
  \item \kw{Perché è interessante}: collega il corso al machine learning e pone subito il
        problema computazionale (con 20 feature, $2^{20}$ coalizioni da valutare).
  \item \kw{Esperimenti}: Shapley esatto su un dataset piccolo, confronto con
        l'approssimazione per campionamento e con la libreria \texttt{shap}; verifica degli
        assiomi (efficienza, giocatore nullo) sui risultati.
  \item \kw{Rischio}: medio. Va definito con cura che cosa vuol dire ``usare solo le feature
        in $C$''.
\end{itemize}

\subsection*{Idea 11: struttura di coalizione per formare squadre \hfill\sbloccata{5}}
\begin{itemize}
  \item \kw{Modello}: assegnare persone a gruppi di lavoro (o veicoli a rotte, o server a
        compiti) quando il valore di un gruppo non è additivo. Non superadditivo: gruppi
        troppo grandi rendono meno.
  \item \kw{Perché è interessante}: la programmazione dinamica esatta in $O(3^n)$ diventa
        impraticabile già con poche decine di agenti; si confronta con euristiche (greedy,
        clustering, simulated annealing), misurando quanto ci si allontana dall'ottimo
        dove questo si può ancora calcolare.
  \item \kw{Rischio}: medio. Aiuta avere già il codice esatto per avere un riferimento.
\end{itemize}

\section{Argomenti ancora da affrontare}

\begin{center}
\renewcommand{\arraystretch}{1.25}
\begin{tabularx}{\linewidth}{@{}l X@{}}
\toprule
\textbf{Argomento} & \textbf{Anticipazioni (possibili spunti)} \\
\midrule
Giochi coalizionali (seconda parte) & il notebook del lab 5 annuncia rappresentazioni compatte
 e approssimazioni; potere di voto, centralità nelle reti \\
Scelta sociale computazionale  & regole di voto, paradosso di Condorcet; \en{ensemble} di modelli \\
Matching markets               & matching stabile; studenti--alloggi, donatori di rene, lavoratori--compiti \\
Mechanism design               & aste, incentivi a dire il vero \\
Fair division                  & divisione equa di risorse \\
\bottomrule
\end{tabularx}
\end{center}

% =====================================================================
\chapter{Domande aperte per il docente}
\label{cap:domande}

Da chiarire il prima possibile, perché cambiano il modo di lavorare.

\begin{enumerate}
  \item Il progetto è \kw{individuale} o \kw{di gruppo}? \hfill\aperta
  \item Il \kw{problema} lo sceglie lo studente o viene assegnato? Serve un'approvazione
        preventiva della proposta? \hfill\aperta
  \item Cosa si \kw{consegna}: codice, relazione scritta, presentazione? Ci sono limiti di
        lunghezza? \hfill\aperta
  \item C'è un \kw{linguaggio} richiesto (Python, come nei lab)? \hfill\aperta
  \item Quando va consegnato rispetto alla data dell'appello? \hfill\aperta
  \item All'esame c'è anche una \kw{discussione orale} o domande di teoria? \hfill\aperta
  \item Quali sono i \kw{criteri di valutazione}? \hfill\aperta
\end{enumerate}

% =====================================================================
\chapter{Piano di lavoro e decisioni}

\section{Calendario indicativo}

Proposta da aggiornare quando saranno note le date ufficiali.

\begin{center}
\renewcommand{\arraystretch}{1.25}
\begin{tabularx}{\linewidth}{@{}l X@{}}
\toprule
\textbf{Periodo} & \textbf{Obiettivo} \\
\midrule
Ottobre    & seguire le lezioni, aggiungere idee alla banca, avere risposta alle domande aperte \\
Novembre   & \textbf{siamo qui}: le domande aperte sono ancora senza risposta e il tema non è
             scelto. Scegliere il tema entro fine mese; scrivere la formalizzazione
             (Capitolo~\ref{cap:metodo}, domande 1--7) \\
Dicembre   & implementazione ed esperimenti \\
Gennaio    & relazione, rifinitura, preparazione della discussione \\
\bottomrule
\end{tabularx}
\end{center}

\section{Registro delle decisioni}

\begin{center}
\renewcommand{\arraystretch}{1.25}
\begin{tabularx}{\linewidth}{@{}l X@{}}
\toprule
\textbf{Data} & \textbf{Decisione} \\
\midrule
29/09/2026 & Creato il documento di progetto. Nessun tema scelto. \\
03/10/2026 & Riorganizzato il repository in \texttt{lessons/}, \texttt{labs/},
             \texttt{documentazione/}. Aggiunta la Lezione~2 e due nuove idee
             (ricerca bayesiana, atteggiamento verso il rischio). Nessun tema scelto. \\
02/11/2026 & Aggiunte le Lezioni 3--5 (giochi non cooperativi, equilibri correlati e
             apprendimento, giochi coalizionali) con i laboratori 3--5, e sei nuove idee
             (6--11). Aggiornata la Lezione 2 con le slide del Lab 2: con le navi affondate
             la ricerca bayesiana scende da 54 a 41--44 colpi. Nessun tema scelto. \\
\bottomrule
\end{tabularx}
\end{center}

\end{document}
